
\documentclass{article}
\usepackage{colortbl}
\usepackage{makecell}
\usepackage{multirow}
\usepackage{supertabular}

\begin{document}

\newcounter{utterance}

\centering \large Interaction Transcript for game `cladder', experiment `full\_v1.5\_default', episode 4040 with Qwen3.5{-}27B.
\vspace{24pt}

{ \footnotesize  \setcounter{utterance}{1}
\setlength{\tabcolsep}{0pt}
\begin{supertabular}{c@{$\;$}|p{.15\linewidth}@{}p{.15\linewidth}p{.15\linewidth}p{.15\linewidth}p{.15\linewidth}p{.15\linewidth}}
    \# & \multicolumn{2}{c}{Player} && \multicolumn{2}{c}{Game Master} \\
    \hline

    \theutterance \stepcounter{utterance}  
    & & & \multicolumn{4}{p{0.6\linewidth}}{
        \cellcolor[rgb]{0.9,0.9,0.9}{
            \makecell[{{p{\linewidth}}}]{
                \texttt{\tiny{[P1$\langle$GM]}}
                \texttt{You are an expert in causal inference. The following question is not a typical commonsense query, but rather a meticulously designed question created by a professor specializing in causal inference, intended to assess the students' mastery of the course content.} \\
\texttt{Imagine a self{-}contained, hypothetical world with only the following conditions, and without any unmentioned factors or causal relationships: Smoking has a direct effect on effort and lactose intolerance. Effort has a direct effect on lactose intolerance. For nonsmokers, the probability of being lactose intolerant is 77\%. For smokers, the probability of being lactose intolerant is 37\%. Will smoking decrease the chance of being lactose intolerant?} \\
\texttt{Start your answer with 'Yes' or 'No', followed by additional reasoning or evidence to support your explanation.} \\
            }
        }
    }
     \\ \\

    \theutterance \stepcounter{utterance}  
    & \multicolumn{4}{p{0.6\linewidth}}{
        \cellcolor[rgb]{0.561,0.737,0.561}{
            \makecell[{{p{\linewidth}}}]{
                \texttt{\tiny{[P1$\rangle$GM]}}
                \texttt{No} \\
\\ 
\texttt{The reasoning relies on distinguishing between the **observed association** (marginal probability) and the **causal effect** (interventional probability) in the presence of a confounder.} \\
\\ 
\texttt{1.  **Analyze the Causal Structure**:} \\
\texttt{    *   The problem states: Smoking \$\textbackslash{}to\$ Effort, Smoking \$\textbackslash{}to\$ Lactose Intolerance, and Effort \$\textbackslash{}to\$ Lactose Intolerance.} \\
\texttt{    *   Here, **Effort** acts as a **mediator** (part of the causal path from Smoking to Lactose Intolerance) and potentially a confounder if we were looking at the effect of Effort, but in this specific setup, Smoking is the root cause.} \\
\texttt{    *   Crucially, the problem states there are "no unmentioned factors." This implies there is no hidden common cause (confounder) between Smoking and Lactose Intolerance other than the paths explicitly described.} \\
\\ 
\texttt{2.  **Analyze the Data**:} \\
\texttt{    *   \$P(\textbackslash{}text\{Lactose Intolerant\} \textbackslash{}mid \textbackslash{}text\{Non{-}Smoker\}) = 0.77\$} \\
\texttt{    *   \$P(\textbackslash{}text\{Lactose Intolerant\} \textbackslash{}mid \textbackslash{}text\{Smoker\}) = 0.37\$} \\
\texttt{    *   The observed data shows a strong negative association: Smokers are much less likely to be lactose intolerant than non{-}smokers.} \\
\\ 
\texttt{3.  **Determine the Causal Effect**:} \\
\texttt{    *   The question asks: "Will smoking **decrease** the chance of being lactose intolerant?" This is asking for the causal effect of the intervention \$do(\textbackslash{}text\{Smoking\})\$.} \\
\texttt{    *   In a world with no unmeasured confounders, the causal effect of Smoking on Lactose Intolerance is the sum of the **direct effect** and the **indirect effect** (through Effort).} \\
\texttt{    *   Mathematically, the total causal effect is \$P(Y \textbackslash{}mid do(X=1)) {-} P(Y \textbackslash{}mid do(X=0))\$.} \\
\texttt{    *   Since there are no back{-}door paths (no confounders) from Smoking to Lactose Intolerance, the observed conditional probability equals the interventional probability:} \\
\texttt{        \$\$P(Y \textbackslash{}mid do(X=1)) = P(Y \textbackslash{}mid X=1) = 0.37\$\$} \\
\texttt{        \$\$P(Y \textbackslash{}mid do(X=0)) = P(Y \textbackslash{}mid X=0) = 0.77\$\$} \\
\texttt{    *   The total causal effect is \$0.37 {-} 0.77 = {-}0.40\$.} \\
\\ 
\texttt{4.  **Addressing Potential "Simpson's Paradox" or Mediation Confusion**:} \\
\texttt{    *   Sometimes students confuse this with scenarios where a variable is a *confounder* (e.g., if} \\
            }
        }
    }
    & & \\ \\

    \theutterance \stepcounter{utterance}  
    & & & \multicolumn{2}{p{0.3\linewidth}}{
        \cellcolor[rgb]{0.9,0.9,0.9}{
            \makecell[{{p{\linewidth}}}]{
                \texttt{\tiny{[GM$|$GM]}}
                \texttt{No} \\
            }
        }
    }
    & & \\ \\

    \theutterance \stepcounter{utterance}  
    & & & \multicolumn{2}{p{0.3\linewidth}}{
        \cellcolor[rgb]{0.9,0.9,0.9}{
            \makecell[{{p{\linewidth}}}]{
                \texttt{\tiny{[GM$|$GM]}}
                \texttt{yes} \\
            }
        }
    }
    & & \\ \\

    \theutterance \stepcounter{utterance}  
    & & & \multicolumn{2}{p{0.3\linewidth}}{
        \cellcolor[rgb]{0.9,0.9,0.9}{
            \makecell[{{p{\linewidth}}}]{
                \texttt{\tiny{[GM$|$GM]}}
                \texttt{game\_result = LOSE} \\
            }
        }
    }
    & & \\ \\

\end{supertabular}
}

\end{document}
